\section{The two design questions}\label{sec:q}
\sectionbrief
{the two questions put to the machine-learning side: must the basis functions of the interpolation be chosen in advance, and how should the outer structure around the analytic core be designed.}
{to show, in one place, what the algorithm now answers and with what evidence, so that the remaining open parts are visible.}
{summary of Sections~\ref{sec:interp}--\ref{sec:loop} from the point of view of the two questions.}

\subsection{Q1: must the basis functions be chosen in advance?}
\textbf{Inside the class of linear constant-coefficient equations in $1+1$ dimensions, no.} The interpolation is a Gaussian process, which needs no basis, and the ``basis'' used by the core is the family of exponentials $e^{\lambda t+i\kappa x}$, which is universal in this class (Section~\ref{sec:prelim}, \thy{T25}). What is chosen in advance is only a \emph{bound}: factors up to degree $m_j\le2$ in $\tau$ (total $\le3$) and $dx_j\le4$ in $\xi$, plus the class itself. A transport equation is found because its curve $\lambda=-i\kappa$ is present among the fitted points, not because a ``transport basis function'' was in a list.

\begin{caution}[title=Where the answer stops]
Outside the class (variable coefficients, nonlinear terms) there is no universal basis, and the dictionary problem returns, with the theorem that no countable dictionary identifies every equation exactly and that with a fixed dictionary the equation is recovered only modulo its differential consequences. The core's method depends on constant coefficients: extending it needs new work (for example frozen-coefficient fits per cell, which the quadtree would support). The kernel family of the GP is also a choice, although its scales are marginalised and the family gets a posterior.
\end{caution}

\subsection{Q2: the outer structure}
\begin{table}[!htb]
\centering\small
\begin{tabular}{@{}L{3.0cm}L{5.2cm}L{4.4cm}L{1.8cm}@{}}
\toprule
Sub-question & Present answer & Evidence & Status\\\midrule
Loop and stopping rule & recurse by Bayes factors; stop when clean with probability $>0.95$ & tests: heat does not grow, transport plus diffusion is grouped by factor, a fifth-order law grows once under the held-out guard; thresholds are hand defaults & \impl{}\\
Parameters from $\mathcal W$ & DP mixture on the cloud gives $P(k)$, centres $\bar{\bs\beta}_j$, covariances $\Sigma_j$ & $P(k=2)=1$ on transport plus diffusion, $P(k=1)$ on heat; merge step thin & \implu{}\\
Encoder and clustering & analytic features; GMM, $k$-means, spectral, agglomerative on rectangles & ARI $1.0$ with exact components, aligned windows; $0.63$ if straddled; dominant balance $0.2$--$0.35$ & \implu{}\\
Uncertainty through a discontinuous map & Monte Carlo over GP samples; never linearised & GP band coverage within $\pm5$ points in $11$ of $12$ cells; the end-to-end benchmark below runs on GP samples & \impl{}\\
Noise and validation & benchmark against PDE-FIND, weak-form, ensemble and Bayesian SINDy, BOP-DMD & full pipeline, run~6 (30 runs per cell, right factor set at noise $10^{-3},10^{-2},5\cdot10^{-2}$): grid $1.00/1.00/0.93$, $2000$ points $1.00/0.97/0.90$, $400$ points $1.00/0.93/0.73$, at least as good as every baseline in all nine cells. After the later fixes, final benchmark (10 cases $\times$ grid $1\%$, grid $5\%$, $400$ points $5\%$): $30/30$ right, no false-confident; best baseline (weak SINDy) $23$ right, $5$ false-confident & \impl{}\\
\bottomrule
\end{tabular}
\caption{The second question, point by point.}
\end{table}

\begin{keypoint}
The point on which the whole proposal depends, ``scattered noisy data $\to$ GP samples $\to$ core $\to$ loop'', is demonstrated on the ten named cases, which also tuned the defaults. Generality is tested on randomly drawn equations (Section~\ref{sec:limits}): $20$ draws, no false-confident answer, but a small sample.
\end{keypoint}
